% Responsable(s): por asignar
% Objetivo: responder la pregunta de investigación, resumir los hallazgos
% principales y señalar limitaciones o líneas de trabajo futuras.

% Redacción pendiente.
Los estudios revisados permiten concluir de manera preliminar que el
rendimiento de un sistema de memoria no depende únicamente de la tecnología
utilizada, sino también de factores como la latencia, el ancho de banda, la
organización de la jerarquía de memoria y la arquitectura del sistema. Los
resultados analizados muestran que tecnologías con un mayor ancho de banda,
como HBM, no necesariamente presentan los menores tiempos de acceso, por lo
que no existe una única alternativa que sea superior para todas las cargas de
trabajo \cite{ibeid2025aurora,esmailidokht2024mess}.

Asimismo, las investigaciones recientes muestran que el rendimiento real puede
diferir considerablemente de los valores teóricos debido al controlador de
memoria, la microarquitectura, la ubicación de los datos y los patrones de
acceso de las aplicaciones \cite{eyerman2022dram,velten2022memory}. Este
comportamiento también se observa en tecnologías de alto ancho de banda, cuyo
rendimiento efectivo depende de la manera en que se organizan y realizan los
accesos a memoria \cite{huang2022shuhai}. Tecnologías como CXL y las memorias
heterogéneas amplían las posibilidades de capacidad y organización del sistema,
aunque también pueden introducir costos adicionales de latencia
\cite{wang2024hitchhiker,lurbe2023controlador}.

A partir de la evidencia recopilada, la investigación resulta viable, ya que
existen estudios recientes con mediciones cuantitativas que permiten comparar
diferentes tecnologías y configuraciones de memoria. De forma preliminar, puede
establecerse que la selección de una tecnología de memoria debe realizarse
considerando conjuntamente el tiempo de acceso, la latencia, el ancho de banda
y las necesidades específicas de cada aplicación. Estos resultados serán
complementados con el análisis de las demás tecnologías incluidas en la
investigación.